\chapter{Ordinale, Hierarchien und Modellkonstruktionen}

\section{Ordinale und Rekursion}

\FormulaDefDeltaK[Ordinalzahl]
{\begin{aligned}
 \alpha\in\BXLVIIIOrd\quad:\Longleftrightarrow\quad&
 \alpha\text{ ist transitiv und durch }\in\text{ wohlgeordnet},\\
 \alpha<\beta\quad:\Longleftrightarrow\quad&\alpha\in\beta
 \qquad(\alpha,\beta\in\BXLVIIIOrd)
\end{aligned}}
{B48OrdinalDef}{
  \DeltaRow{Transitivität}{\forall x\in\alpha\ (x\subseteq\alpha)}
}
\noindent
Es gelten \(0=\varnothing\), \(\alpha+1=\alpha\cup\{\alpha\}\);
eine von \(0\) verschiedene Ordinalzahl ohne unmittelbaren Vorgänger
heißt Limesordinalzahl. \(\BXLVIIIOrd\) bezeichnet eine definierbare
Klasse, nicht die Menge aller Ordinalzahlen.

\FormulaThmDeltaK[Vergleich von Ordinalzahlen]
{\alpha,\beta\in\BXLVIIIOrd\ \vdash\
 (\alpha\in\beta)\ \dot\lor\ (\alpha=\beta)\ \dot\lor\ (\beta\in\alpha)}
{B48OrdinalComparison}{
  \DeltaRow{Ordinale}{\alpha,\beta}
}
\begin{tabproof}
 \proofstep{1}{\alpha,\beta\in\BXLVIIIOrd}{\rA}
 \proofstep{1}{\BXLVIIIText{Jedes Element einer Ordinalzahl ist selbst
 transitiv und durch \(\in\) wohlgeordnet: Transitivität folgt aus der
 Transitivität der Ordnungsrelation, und die Wohlordnung wird vererbt.}}
 {\BXLVIIIWhy{Definition einer Ordinalzahl}}
 \proofstep{1}{\BXLVIIIText{Eine transitive echte Teilmenge \(\gamma\)
 einer Ordinalzahl \(\alpha\) ist ein Anfangsstück.
 Ist \(\delta\) das kleinste Element von \(\alpha\setminus\gamma\),
 so ist \(\gamma=\delta\): Alle kleineren Elemente liegen in
 \(\gamma\), ein größeres Element würde wegen Transitivität auch
 \(\delta\) enthalten.}}
 {\BXLVIIIWhy{Wohlordnung und Transitivität}}
 \proofstep{1}{\BXLVIIIText{Der Schnitt \(\gamma=\alpha\cap\beta\)
 ist transitiv. Wäre er in beiden Ordinalzahlen echt, so wäre nach
 der vorigen Zeile \(\gamma\in\alpha\cap\beta=\gamma\), ein
 Widerspruch. Also ist eine der Ordinalzahlen Teilmenge der anderen.}}
 {\BXLVIIIWhy{Fundierung}}
 \proofstep{1}{\BXLVIIIText{Bei echter Inklusion liefert die
 Anfangsstückbeschreibung \(\alpha\in\beta\) beziehungsweise
 \(\beta\in\alpha\). Gleichheit ist der verbleibende Fall.
 Die Fälle schließen einander durch Fundierung aus.}}
 {\BXLVIIIWhy{Anfangsstücke}}
\end{tabproof}

\FormulaThmDeltaK[Transfinite Induktion]
{\bigl[\forall\alpha\in\BXLVIIIOrd\
 ((\forall\beta<\alpha\ \psi(\beta))\to\psi(\alpha))\bigr]
 \ \vdash\ \forall\alpha\in\BXLVIIIOrd\ \psi(\alpha)}
{B48TransfiniteInduction}{
  \DeltaRow{Formel}{\psi(\alpha)\text{ mit festen Parametern}}
}
\begin{tabproof}
 \proofstep{1}{\forall\alpha\ ((\forall\beta<\alpha\ \psi(\beta))\to\psi(\alpha))}
 {\rA}
 \proofstep{2}{\exists\gamma\in\BXLVIIIOrd\ \neg\psi(\gamma)}{\rA}
 \proofstep{2}{\BXLVIIIText{Wähle ein solches \(\gamma\). Die nichtleere
 Menge \(\{\xi\le\gamma:\neg\psi(\xi)\}\) besitzt ein kleinstes
 Element \(\delta\). Für jedes \(\xi<\delta\) gilt \(\psi(\xi)\).}}
 {\BXLVIIIWhy{Aussonderung in \(\gamma+1\), Wohlordnung}}
 \proofstep{1,2}{\psi(\delta)\land\neg\psi(\delta)}
 {\BXLVIIIWhy{Induktionsvoraussetzung an \(\delta\)}}
 \proofstep{1}{\forall\alpha\in\BXLVIIIOrd\ \psi(\alpha)}
 {\BXLVIIIWhy{Widerspruchsbeweis}}
\end{tabproof}

\FormulaThmDeltaK[Transfinite Rekursion]
{\begin{gathered}
 H\text{ ordnet jeder ordinal indizierten Mengenfolge genau eine Menge zu}\\
 \vdash\ \exists! F\quad
 \forall\alpha\in\BXLVIIIOrd\ \bigl(F(\alpha)=H(F\mathbin{\upharpoonright}\alpha)\bigr)
 \end{gathered}}
{B48TransfiniteRecursion}{
  \DeltaRow{Klassenfunktion}{H\text{ durch eine Formel definiert}}
}
\begin{tabproof}
 \proofstep{1}{\BXLVIIIText{\(H\) ist auf allen ordinal indizierten
 Mengenfolgen eindeutig und überall definiert.}}{\rA}
 \proofstep{1}{\BXLVIIIText{Eine Approximation ist eine Funktion mit
 ordinalem Definitionsbereich, welche die Rekursionsgleichung an jeder
 Stelle ihres Bereichs erfüllt. Zwei Approximationen stimmen auf dem
 gemeinsamen Bereich überein: Ein kleinster Unterschied hätte gleiche
 Vorgängerfolgen und damit denselben durch \(H\) bestimmten Wert.}}
 {\BXLVIIIWhy{Transfinite Induktion}}
 \proofstep{1}{\BXLVIIIText{Die leere Funktion ist eine Approximation.
 Aus einer Approximation \(f\) auf \(\alpha\) entsteht eine auf
 \(\alpha+1\), indem das Paar \((\alpha,H(f))\) angefügt wird.}}
 {\BXLVIIIWhy{Leere Menge, Paarung und Vereinigung}}
 \proofstep{1}{\BXLVIIIText{An einer Limesstelle \(\lambda\) sammelt
 Ersetzung die eindeutig bestimmten Approximationen auf
 \(\alpha<\lambda\). Ihre Vereinigung ist wegen der Übereinstimmung
 wieder eine Approximation.}}
 {\BXLVIIIWhy{Ersetzung und Vereinigung}}
 \proofstep{1}{\BXLVIIIText{Transfinite Induktion liefert jede
 Mengenapproximation. Die Formel, welche ihre übereinstimmenden Werte
 beschreibt, definiert \(F\). Die erste Zeile des Arguments gibt
 zugleich die Eindeutigkeit.}}
 {\BXLVIIIWhy{Keine Vereinigung zu einer Menge über alle Ordinale nötig}}
\end{tabproof}

\FormulaThmDeltaK[Rekursion auf einer fundierten Mengenrelation]
{\begin{gathered}
 R\subseteq X^2\text{ fundiert},\ H\text{ eindeutig und überall definiert}\\
 \vdash\ \exists! f:X\longrightarrow V\quad
 f(x)=H\bigl(x,f\mathbin{\upharpoonright}\{y:yRx\}\bigr)
 \end{gathered}}
{B48WellFoundedRecursion}{
  \DeltaRow{Menge}{X}
  \DeltaRow{Relation}{R}
}
\begin{tabproof}
 \proofstep{1}{R\subseteq X^2\text{ fundiert}}{\rA}
 \proofstep{1}{\BXLVIIIText{Auf vorgängerabgeschlossenen Teilmengen
 definierte Lösungen stimmen auf Überschneidungen überein.
 Andernfalls besitzt die Menge ihrer Unterschiede ein
 \(R\)-minimales Element; dort stimmen die Vorgängerwerte schon
 überein.}}
 {\BXLVIIIWhy{Fundiertheit}}
 \proofstep{1}{\BXLVIIIText{Nenne \(x\) erreichbar, wenn eine Lösung
 auf der Vorgängerhülle von \(x\), einschließlich \(x\), existiert.
 Sind alle unmittelbaren Vorgänger erreichbar, so sammelt Ersetzung
 deren eindeutige Lösungen. Ihre Vereinigung liefert die Werte unter
 \(x\), und \(H\) liefert den Wert bei \(x\).}}
 {\BXLVIIIWhy{Vorgängerhülle durch endliche \(R\)-Ketten;
 Ersetzung}}
 \proofstep{1}{\BXLVIIIText{Die Menge nicht erreichbarer Elemente
 könnte kein minimales Element haben; sie ist daher leer.
 Ersetzung und Vereinigung sammeln die kompatiblen lokalen Lösungen
 zur gesuchten Funktion \(f\).}}
 {\BXLVIIIWhy{Fundiertheit und Aussonderung}}
\end{tabproof}

\FormulaThmDeltaK[Mostowski-Kollaps]
{\begin{gathered}
 R\subseteq X^2\text{ fundiert und extensional}\\
 \vdash\ \exists! (N,\pi)\quad
 N\text{ transitiv},\quad \pi:(X,R)\cong(N,\in)
 \end{gathered}}
{B48Mostowski}{
  \DeltaRow{Extensionalität}{
  \forall x,x'\in X\
  (\{y:yRx\}=\{y:yRx'\}\to x=x')}
}
\begin{tabproof}
 \proofstep{1}{R\text{ fundiert und extensional auf }X}{\rA}
 \proofstep{1}{\pi(x)=\{\pi(y):yRx\}}
 {\FormulaRefAuto{B48WellFoundedRecursion}[thm]}
 \proofstep{1}{\BXLVIIIText{Wäre \(\pi\) nicht injektiv, wähle ein
 \(R\)-minimales \(x\), das an einer Kollision
 \(\pi(x)=\pi(x')\), \(x\ne x'\), beteiligt ist. Jeder Vorgänger
 \(yRx\) besitzt dann ein eindeutiges Urbild seines Bildes.
 Die Gleichheit der Bildmengen erzwingt deshalb in beiden Richtungen
 \(\{y:yRx\}=\{y:yRx'\}\).}}
 {\BXLVIIIWhy{Fundiertheit; Rekursionsgleichung}}
 \proofstep{1}{\BXLVIIIText{Extensionalität gäbe \(x=x'\), also ist
 \(\pi\) injektiv. Ihre Bildmenge \(N=\pi[X]\) ist nach der
 Rekursionsgleichung transitiv, und \(yRx\) gilt genau dann, wenn
 \(\pi(y)\in\pi(x)\).}}
 {\BXLVIIIWhy{Ersetzung und Injektivität}}
 \proofstep{1}{\BXLVIIIText{Jeder andere Isomorphismus auf eine
 transitive Menge erfüllt dieselbe Rekursionsgleichung.
 Die Eindeutigkeit der Rekursion liefert \(\pi\) und damit \(N\).}}
 {\BXLVIIIWhy{Fundierte Rekursion}}
\end{tabproof}

\noindent
Eine strikte Wohlordnung ist extensional und fundiert. Ihr Kollaps ist
eine Ordinalzahl, ihr eindeutig bestimmter \emph{Ordnungstyp}.

\FormulaThmDeltaK[Ranghierarchie]
{\begin{aligned}
 V_0&=\varnothing,& V_{\alpha+1}&=\mathcal P(V_\alpha),&
 V_\lambda&=\bigcup_{\alpha<\lambda}V_\alpha,\\
 \operatorname{rk}(x)&=\sup\{\operatorname{rk}(y)+1:y\in x\},&
 x&\in V_{\operatorname{rk}(x)+1}
\end{aligned}}
{B48RankHierarchy}{
  \DeltaRow{Limesordinal}{\lambda}
  \DeltaRow{Menge}{x}
}
\begin{tabproof}
 \proofstep{}{\BXLVIIIText{Die drei Rekursionsklauseln definieren
 die Mengen \(V_\alpha\). Induktion zeigt ihre Transitivität und
 \(V_\alpha\subseteq V_\beta\) für \(\alpha\le\beta\).}}
 {\FormulaRefAuto{B48TransfiniteRecursion}[thm]}
 \proofstep{}{\BXLVIIIText{Die Vereinigung der durch wiederholte
 Vereinigung aus \(\{x\}\) gewonnenen Mengen ist eine transitive
 Menge, welche \(x\) enthält. Auf ihr ist \(\in\) fundiert.}}
 {\BXLVIIIWhy{Rekursion auf \(\omega\), Fundierungsaxiom}}
 \proofstep{}{\BXLVIIIText{Fundierte Rekursion mit dem Supremum
 der Nachfolger der Vorgängerwerte definiert den Rang.
 Alle Werte sind Ordinalzahlen, da Vereinigungen von Mengen von
 Ordinalzahlen wieder Ordinalzahlen sind.}}
 {\FormulaRefAuto{B48WellFoundedRecursion}[thm]}
 \proofstep{}{\BXLVIIIText{Induktion entlang \(\in\) gibt
 \(y\in V_{\operatorname{rk}(y)+1}\) und
 \(\operatorname{rk}(y)+1\le\operatorname{rk}(x)\) für \(y\in x\).
 Also \(x\subseteq V_{\operatorname{rk}(x)}\), wie behauptet.}}
 {\BXLVIIIWhy{Potenzmengenstufe und Monotonie}}
\end{tabproof}

\section{Kardinale und die benötigte Arithmetik}

\FormulaDefDeltaK[Kardinalzahl und Vergleich]
{\begin{aligned}
 |A|&=\min\{\alpha\in\BXLVIIIOrd:A\approx\alpha\},\\
 \kappa\text{ Kardinalzahl}&\ \Longleftrightarrow\
 \kappa\in\BXLVIIIOrd\land|\kappa|=\kappa,\\
 |A|\le|B|&\ \Longleftrightarrow\
 \text{es gibt eine Injektion }A\longrightarrow B
\end{aligned}}
{B48CardinalDef}{
  \DeltaRow{Grundlage}{\mathsf{AC}\text{ und der Wohlordnungssatz}}
}
\noindent
Der Wohlordnungssatz aus den bisherigen Bänden liefert einen
Ordinalzahlvertreter. Unterhalb eines gewählten Vertreters kann das
Minimum durch Aussonderung gebildet werden. Damit ist die Definition
keine Auswahl aus einer echten Klasse ohne Begründung.

\FormulaThmDeltaK[Hartogs und Nachfolgerkardinale]
{\forall\kappa\text{ Kardinalzahl}\quad
 \exists\kappa^+\quad
 \kappa^+=\min\{\lambda\text{ Kardinalzahl}:\kappa<\lambda\}}
{B48Hartogs}{
  \DeltaRow{Kardinalzahl}{\kappa}
}
\begin{tabproof}
 \proofstep{1}{\kappa\text{ Kardinalzahl}}{\rA}
 \proofstep{1}{\BXLVIIIText{Die Paare \((D,R)\) mit \(D\subseteq\kappa\)
 und \(R\subseteq D^2\) bilden eine Menge. Sondere die Wohlordnungen
 aus und sammle mit Ersetzung ihre eindeutigen Ordnungstypen zu \(S\).}}
 {\BXLVIIIWhy{Potenzmengen, Aussonderung, Mostowski-Kollaps}}
 \proofstep{1}{\BXLVIIIText{Die Ordinalzahl
 \(h=\sup\{\alpha+1:\alpha\in S\}\) injiziert nicht in \(\kappa\).
 Eine solche Injektion würde eine Wohlordnung auf ihrem Bild mit
 Typ \(h\) liefern; dann wäre \(h\in S\) und \(h+1\le h\).}}
 {\BXLVIIIWhy{Widerspruch}}
 \proofstep{1}{\BXLVIIIText{Die kleinste Ordinalzahl ohne Injektion
 in \(\kappa\) ist eine Kardinalzahl: Eine Bijektion mit einer
 kleineren Ordinalzahl würde eine solche Injektion ergeben.
 Sie ist größer als \(\kappa\) und kleiner oder gleich jeder
 größeren Kardinalzahl.}}
 {\BXLVIIIWhy{Minimalität; Ordinalvergleich}}
\end{tabproof}
\noindent
Wir setzen \(\aleph_0=\omega\), \(\aleph_1=\aleph_0^+\) und
\(\aleph_2=\aleph_1^+\). \(\omega\) ist die kleinste unendliche
Ordinalzahl; alle ihre Vorgänger sind die natürlichen Zahlen.

\FormulaThmDeltaK[Cantor--Schröder--Bernstein]
{A\hookrightarrow B,\ B\hookrightarrow A\ \vdash\ A\approx B}
{B48CSB}{
  \DeltaRow{Injektionen}{f:A\longrightarrow B,\quad g:B\longrightarrow A}
}
\begin{tabproof}
 \proofstep{1}{f,g\text{ injektiv}}{\rA}
 \proofstep{1}{\begin{aligned}
 A_0&=A\setminus g[B],\\
 A_{n+1}&=g[f[A_n]],\\
 C&=\bigcup_{n<\omega}A_n
 \end{aligned}}
 {\BXLVIIIWhy{Rekursion und Vereinigung}}
 \proofstep{1}{g(y)\in C\ \Longleftrightarrow\ y\in f[C]\quad(y\in B)}
 {\BXLVIIIWhy{\(g(y)\notin A_0\); Rekursion und Injektivität von \(g\)}}
 \proofstep{1}{h(x)=
 \begin{cases}f(x),&x\in C,\\ g^{-1}(x),&x\in A\setminus C\end{cases}}
 {\BXLVIIIWhy{\(A\setminus C\subseteq g[B]\)}}
 \proofstep{1}{\BXLVIIIText{Der erste Teil ist eine Bijektion von
 \(C\) auf \(f[C]\), der zweite nach Zeile 3 eine Bijektion von
 \(A\setminus C\) auf \(B\setminus f[C]\).
 Die disjunkten Teile ergeben die Bijektion \(h:A\to B\).}}
 {\BXLVIIIWhy{Zusammensetzen von Bijektionen}}
\end{tabproof}

\FormulaThmDeltaK[Produkte unendlicher Kardinale]
{\kappa\ge\aleph_0\ \vdash\
 \kappa\cdot\kappa=\kappa,\quad
 \kappa\cdot\aleph_0=\kappa,\quad
 |\,\kappa^{<\omega}\,|=\kappa}
{B48CardinalProduct}{
  \DeltaRow{Kardinalzahl}{\kappa}
}
\begin{tabproof}
 \proofstep{1}{\kappa\text{ unendliche Kardinalzahl}}{\rA}
 \proofstep{1}{\BXLVIIIText{Induktion über die unendlichen Kardinale.
 Ordne \(\kappa^2\) zuerst nach
 \(\max(\alpha,\beta)\), bei Gleichheit lexikographisch.
 Dies ist eine Wohlordnung.}}
 {\BXLVIIIWhy{Ordinalwohlordnung und lexikographischer Vergleich}}
 \proofstep{1}{\BXLVIIIText{Das Vorgängerstück eines Paares mit
 Maximum \(\gamma\) liegt in \((\gamma+1)^2\).
 Dessen Mächtigkeit ist kleiner als \(\kappa\):
 im endlichen Fall unmittelbar, sonst nach der
 Induktionsvoraussetzung für \(|\gamma+1|<\kappa\).}}
 {\BXLVIIIWhy{Induktionsschritt; bei \(\kappa=\omega\) nur endliche Stücke}}
 \proofstep{1}{\BXLVIIIText{Der Ordnungstyp kann nicht größer als
 \(\kappa\) sein, sonst hätte das Element an der Stelle \(\kappa\)
 ein Vorgängerstück der Mächtigkeit \(\kappa\).
 Also \(|\kappa^2|\le\kappa\). Die Diagonale liefert die
 umgekehrte Ungleichung.}}
 {\BXLVIIIWhy{Ordnungstyp und Cantor--Schröder--Bernstein}}
 \proofstep{1}{\BXLVIIIText{Da \(\omega\le\kappa\), folgt
 \(\kappa\cdot\omega=\kappa\). Endliche Potenzen von \(\kappa\)
 haben nach Induktion ebenfalls Mächtigkeit \(\kappa\);
 ihre abzählbare Vereinigung hat daher Mächtigkeit \(\kappa\).}}
 {\BXLVIIIWhy{Produktgleichung; Auswahl von Injektionen}}
\end{tabproof}
\noindent
Insbesondere hat eine Vereinigung von höchstens \(\mu\) Mengen mit
je höchstens \(\nu\) Elementen höchstens \(\mu\cdot\nu\) Elemente:
Wähle die Einzelinjektionen und codiere jedes Element mit dem kleinsten
Index, bei dem es auftritt, und seinem Bild. Hier wird Auswahl benutzt.

\FormulaThmDeltaK[Potenz- und Exponentengesetze]
{\begin{aligned}
 |\mathcal P(A)|&=2^{|A|},\\
 (\kappa^\lambda)^\mu&=\kappa^{\lambda\cdot\mu},\\
 \lambda\le\mu&\ \Longrightarrow\ 2^\lambda\le2^\mu
\end{aligned}}
{B48ExponentLaws}{
  \DeltaRow{Kardinale}{\kappa,\lambda,\mu}
  \DeltaRow{Definition}{\kappa^\lambda=|\{f:f:\lambda\longrightarrow\kappa\}|}
}
\begin{tabproof}
 \proofstep{}{\BXLVIIIText{Die Zuordnung \(S\mapsto\chi_S\) ist
 eine Bijektion von \(\mathcal P(A)\) auf die Funktionen von
 \(A\) nach \(2=\{0,1\}\).}}
 {\BXLVIIIWhy{Charakteristische Funktionen}}
 \proofstep{}{\BXLVIIIText{Zu \(F:\mu\to(\lambda\to\kappa)\) gehört
 \(f:\mu\times\lambda\to\kappa\) mit \(f(i,j)=F(i)(j)\).
 Die umgekehrte Zuordnung \(f\mapsto(i\mapsto(j\mapsto f(i,j)))\)
 ist invers.}}
 {\BXLVIIIWhy{Explizite Bijektion}}
 \proofstep{}{\BXLVIIIText{Aus einer Injektion
 \(i:\lambda\to\mu\) entsteht die Injektion
 \(S\mapsto i[S]\) von \(\mathcal P(\lambda)\) nach
 \(\mathcal P(\mu)\).}}
 {\BXLVIIIWhy{Injektivität der Bildzuordnung}}
\end{tabproof}

\section{Elementare Teilstrukturen}

Eine Teilstruktur \(X\) einer Struktur \(A\) heißt elementar,
geschrieben \(X\prec A\), wenn jede Formel mit Parametern aus \(X\)
in beiden Strukturen denselben Wahrheitswert besitzt.

\FormulaThmDeltaK[Abzählbare Skolemhülle]
{\begin{gathered}
 S\subseteq A,\quad \mathcal L\text{ abzählbar}\\
 \vdash\ \exists X\prec A\quad
 S\subseteq X,\qquad |X|\le\max(|S|,\aleph_0)
 \end{gathered}}
{B48SkolemHull}{
  \DeltaRow{Mengenstruktur}{A}
  \DeltaRow{Parametermenge}{S}
}
\begin{tabproof}
 \proofstep{1}{S\subseteq A,\quad\mathcal L\text{ abzählbar}}{\rA}
 \proofstep{1}{\BXLVIIIText{Wohlordne die Grundmenge von \(A\).
 Wähle für jede Existenzformel und jedes passende endliche
 Parametertupel den kleinsten Zeugen, falls einer existiert.
 Füge gegebenenfalls ein festes Element als Startwert hinzu.}}
 {\BXLVIIIWhy{Auswahl; Erfüllungsrelation für Mengenstrukturen}}
 \proofstep{1}{\BXLVIIIText{Beginne mit \(S\) und schließe in
 abzählbar vielen Schritten unter diesen Zeugenfunktionen ab.
 Die Vereinigung \(X\) hat Mächtigkeit höchstens
 \(\max(|S|,\aleph_0)\), weil es nur abzählbar viele Funktionen
 endlicher Stelligkeit gibt.}}
 {\FormulaRefAuto{B48CardinalProduct}[thm]}
 \proofstep{1}{\BXLVIIIText{Formelinduktion beweist \(X\prec A\).
 Atomare Formeln stimmen wegen der Teilstruktur überein,
 Junktoren nach Induktion. Ist eine Existenzformel in \(A\)
 mit Parametern aus \(X\) wahr, so liegen alle Parameter in
 einer endlichen Stufe und ihr gewählter Zeuge in der nächsten.
 Die Rückrichtung benutzt denselben Zeugen in \(A\).}}
 {\BXLVIIIWhy{Zeugenabschluss und Erfüllung}}
\end{tabproof}
